\bookchapter{Глава 68. Путь дальше — разрешения, уведомления и современная архитектура}

\chapterintro{Уметь собрать приложение — только начало. В этой главе — несколько бродов на пути от «умею делать» к «делаю хорошо»: как изящно запрашивать разрешения, как слать уведомления, не докучая, как поднять архитектуру ещё на этаж и — как в 2026 году писать код для iOS в паре с ИИ. Каждая из этих тем встретится в реальном проекте незамедлительно.}

\booksection{Разрешения: сперва стучим, потом входим}

Камера, геолокация, контакты, уведомления… в этих «чувствительных зонах» iOS требует, чтобы приложение при использовании показывало пользователю запрос. Правило — три шага: «объявить + запросить + изящно деградировать»:

\begin{lstlisting}[language=Swift]
// 1. Объявите текст о цели использования в Info.plist (без него — мгновенный краш!)
// Privacy - Location When In Use Usage Description: "Используется для запроса погоды в вашем городе"

// 2. Запрашиваем при использовании
import CoreLocation

@Observable
class LocationModel {
    private let manager = CLLocationManager()

    func request() {
        switch manager.authorizationStatus {
        case .notDetermined: manager.requestWhenInUseAuthorization()
        case .denied, .restricted: /* направить в настройки */
        default: /* разрешено, действуем */
        }
    }
}

// 3. После отказа — мягкая деградация: подсказка + переход в настройки
// UIApplication.shared.open(URL(string: UIApplication.openSettingsURLString)!)
\end{lstlisting}

\begin{tcbwarn}{Внимание: разрешения iOS «одноразовые»}
Окно запроса разрешения в iOS показывается лишь однажды — после отказа пользователь его больше не увидит. Поэтому: выбирайте для запроса момент, «когда пользователь вот-вот собрался этим пользоваться» (а не пачку окон при запуске), и объясняйте в тексте, зачем это вам. Стандарт де-факто — подсказка после отказа («включите геолокацию в Настройках»): не оставляйте пользователя в тупике.

\end{tcbwarn}

\booksection{Уведомления: в правильное время говорить полезное}

Уведомления — канал, которым приложение достаёт пользователя, и функция, которую пользователи «заносят в чёрный список» легче всего. В iOS фреймворк UserNotifications и три дисциплины отправки: только полезное, позволяйте отключить, не докучайте ночью. Базовое применение:

\begin{lstlisting}[language=Swift]
import UserNotifications

// 1. Запрашиваем разрешение на уведомления (один диалог)
UNUserNotificationCenter.current()
    .requestAuthorization(options: [.alert, .sound, .badge]) { granted, _ in
        print(granted ? "разрешено" : "отказано")
    }

// 2. Планируем «ежедневное напоминание о записях в 20:00»
func scheduleDailyReminder() {
    let content = UNMutableNotificationContent()
    content.title = "Напоминание о записях"
    content.body = "Сегодня ещё нет записей"
    content.sound = .default

    var comps = DateComponents()
    comps.hour = 20
    let trigger = UNCalendarNotificationTrigger(dateMatching: comps, repeats: true)

    let request = UNNotificationRequest(
        identifier: "daily-reminder",
        content: content,
        trigger: trigger
    )
    UNUserNotificationCenter.current().add(request)
}
\end{lstlisting}

Внимание: разрешение на уведомления в iOS тоже «показывается один раз», и пользователь может отключить его в системных настройках в любой момент — потому так важна ценность содержимого уведомлений: один раз рассылаете что попало — можете потерять пользователя навсегда.

\booksection{Архитектура: MVVM глубже и тестируемость}

MVVM мы трогали в главе 65 — здесь добавим три момента продвинутого уровня:

\begin{itemize}
  \item Абстракция через протоколы: пусть ViewModel зависит от протокола (protocol), а не от конкретного класса — тогда в тестах можно подставить «фиктивный источник данных»; это краеугольный камень тестируемости;
  \item Единственная ответственность: один ViewModel ведает одним делом (учёт не лезет в погоду); функций стало больше — дробите;
  \item Не переусложняйте: небольшому проекту на одного человека хватает «класса @Observable + абстракция протоколами». Тяжеловесные VIPER и Clean Architecture подождут, пока к вам не постучатся команда и сложность — инструменты служат цели, а не понтам.
\end{itemize}

\begin{tcbinfo}{Совет: «дивиденды экосистемы» SwiftUI}
В экосистеме iOS есть ещё два «больше нигде не водящихся» приёма: виджеты (WidgetKit) позволяют пользователю видеть сведения, не открывая приложение («расходы за месяц» из книги учёта — прямо на рабочем столе), а быстрые команды (App Intents) позволяют вызывать ваши функции через Siri или автоматизации. Оба поддержаны официальными фреймворками, порог входа невысок, а прибавка к пользовательскому опыту огромна — на пути прокачки их стоит разблокировать первыми.

\end{tcbinfo}

\booksection{Новое в 2026 году: разработка iOS с помощью ИИ}

В 2026 году ИИ уже глубоко врос в процесс разработки под Apple. Если хотите ускориться, эта связка вполне обкатана:

\begin{center}
\begin{tabularx}{\linewidth}{|X|X|X|}
\hline
\textbf{Этап} & \textbf{Чем помогает ИИ} & \textbf{Ваша ответственность} \tabularnewline
\hline
Писать код & Встроенные в Xcode подсказки кода + генерация сторонними ИИ-плагинами & Ясно описать требования, проверять то, что получилось \tabularnewline
Искать материалы & Отдавать ИИ на разбор официальную документацию и тексты ошибок & Ключевые выводы перепроверять по официальной документации \tabularnewline
Писать тесты & Просить ИИ сгенерировать кейсы XCTest трёх классов: «нормальные / граничные / ошибочные» & Убедиться, что кейсы соответствуют смыслу задачи \tabularnewline
Чинить баги & Приложить журнал падения + относящийся код — ИИ найдёт и починит & Разобраться в правках и прогнать тесты \tabularnewline
\hline
\end{tabularx}
\end{center}

\begin{center}\small\color{gray}Четыре сценария помощи ИИ в разработке iOS\end{center}

Конкретная методология — как давать ИИ задания, как принимать работу, каких граблей избегать — изложена в 11 главах части 7 этой книги, «Практика программирования с ИИ» (главы 37–47); очень советую читать в связке: синтаксис Swift — из этой части, искусство сотрудничества с ИИ — из седьмой.

\begin{tcbpractice}{Практика: добавьте книге учёта «комплект из трёх улучшений»}
\booksection{Ключевое о конкурентности Swift: @MainActor и Task}

В системе конкурентности Swift 6 два правила, которые новичку необходимо понять:

\begin{itemize}
  \item Интерфейс обновляется только в главном потоке: представления SwiftUI по умолчанию @MainActor, сетевые колбэки возвращаются в главный поток и лишь затем меняют состояние (await MainActor.run \{ ... \} либо пометьте ViewModel как @MainActor);
  \item Асинхронные задачи заводим через Task: модификатор .task \{ \} стартует сам при появлении представления и отменяется при исчезновении — это штатный способ запуска асинхронной работы в SwiftUI; не разбрасывайте голые Task по body.
\end{itemize}

\begin{lstlisting}[language=Swift]
struct WeatherScreen: View {
    @State private var vm = WeatherViewModel()

    var body: some View {
        Group {
            switch vm.state {
            case .loading: ProgressView()
            case .success(let d): WeatherContent(data: d)
            case .error(let m): Text(m)
            }
        }
        .task {                       // начинается при появлении, отменяется при исчезновении
            await vm.search(city: "Пекин")
        }
    }
}
\end{lstlisting}

За обоими правилами стоит дизайн Swift 6 «безопасность от гонок данных» — компилятор укажет на ошибку сразу, а не дождётся падения приложения. Ошибки компилятора — не придирки: они принимают пули за вас.

(1) Добавьте книге учёта «напоминание о записи каждый день в 20:00» (разрешение на уведомления + локальные уведомления); (2) переведите источник данных ViewModel на абстракцию через протокол и напишите модульный тест с подстановкой фиктивных данных; (3) добавьте на рабочий стол виджет «Расходы за месяц». Когда все три улучшения готовы, проект приобретает по-настоящему профессиональный вид.

\end{tcbpractice}

\begin{tcbsummary}{Краткие итоги главы}
\begin{itemize}
  \item Три шага разрешений: объявить в Info.plist (пропустите — падение гарантировано), запросить в момент использования, после отказа — подсказать дорогу в Настройки
  \item Окно разрешения iOS показывается один раз: момент — «вот-вот будем пользоваться», в тексте ясная цель
  \item Три дисциплины уведомлений: только полезное, можно отключить, не докучать; фреймворк UserNotifications
  \item MVVM продвинуто: абстракция протоколами (тестируемость), единственная ответственность, без переусложнения
  \item Дивиденды экосистемы: виджеты WidgetKit + быстрые команды App Intents
  \item Четыре сценария помощи ИИ: писать код, искать материалы, писать тесты, чинить баги; методология — в части 7 (главы 37–47)
\end{itemize}

\end{tcbsummary}
